\chapter{第2章\texorpdfstring{\hspace{0.5em}}{ }质点动力学}

运动是物质本身的属性，但物体如何运动则决定于它们之间的相互作用。研究物体之间的相互作用，以及由于这种相互作用所引起的物体运动状态变化的规律，是动力学的任务。本章在物体运动速率远小于光速的前提下，从动量、角动量、能量三个基本守恒量出发，研究物质运动与物体间相互作用的内在联系，阐明三条守恒定律各自的成立条件和应用范畴。

\section{动量\texorpdfstring{\hspace{0.5em}}{ }动量守恒定律}
    \subsection{动量\texorpdfstring{\hspace{0.5em}}{ }质量}
    质点的\textbf{动量}（momentum）定义为质量与速度的乘积：
    \keyformula{\boldsymbol{p}=m\boldsymbol{v}}
        动量是矢量，方向与速度方向相同；单位 $\mathrm{kg\cdot m\cdot s^{-1}}$，量纲 $\mathrm{MLT^{-1}}$。质点系的总动量是各质点动量的矢量和：
    \[
    \boldsymbol{p}=\sum_{i}\boldsymbol{p}_{i}=\sum_{i}m_{i}\boldsymbol{v}_{i}
    \]

    \textbf{质量}是物体的固有属性，有两种量度：\textbf{惯性质量}（物体惯性大小的量度，$m=\dfrac{F}{a}$）与\textbf{引力质量}（物体间产生引力作用“能力”的量度）。实验表明二者数值成正比、与物体成分和结构无关，选用适当单位后可用同一数值表征，故一般情况下不再区分。在 $v\ll c$ 时，质量不随速度改变。

    \subsection{动量的时间变化率\texorpdfstring{\hspace{0.5em}}{ }力}
        \subsubsection{质点动量的时间变化率\texorpdfstring{\hspace{0.5em}}{ }牛顿第二定律}
        当质点运动状态改变时，动量随时间变化。在 $v\ll c$、质量可视为常量的条件下
        \keyformula{\boldsymbol{F}=\frac{\mathrm{d}\boldsymbol{p}}{\mathrm{d}t}=\frac{\mathrm{d}(m\boldsymbol{v})}{\mathrm{d}t}=m\frac{\mathrm{d}\boldsymbol{v}}{\mathrm{d}t}=m\boldsymbol{a}}
        这就是\textbf{牛顿第二定律}，也是力的定义式：力是质点动量随时间变化率的量度。它是瞬时矢量关系，$F$ 与 $a$ 同时存在、同时消失、方向相同。直角坐标系与自然坐标系中的分量式为
        \[
        F_{x}=ma_{x},\quad F_{y}=ma_{y},\quad F_{z}=ma_{z};
        \qquad F_{\mathrm{t}}=m\frac{\mathrm{d}v}{\mathrm{d}t},\quad F_{\mathrm{n}}=m\frac{v^{2}}{\rho}
        \]
        力的单位是牛顿：$1\ \mathrm{N}=1\ \mathrm{kg\cdot m\cdot s^{-2}}$。

        \subsubsection{牛顿三定律与惯性系}
        \begin{enumerate}
            \item \textbf{牛顿第一定律（惯性定律）}：任何物体都要保持其静止或匀速直线运动状态，直到外力迫使它改变运动状态为止。它指出物体具有保持运动状态的属性（\textbf{惯性}），并定义了\textbf{惯性系}——第一定律成立的参考系；第一定律不成立的参考系称为\textbf{非惯性系}。
            \item \textbf{牛顿第二定律}：$\boldsymbol{F}=\mathrm{d}\boldsymbol{p}/\mathrm{d}t$，只适用于质点，且只在惯性系中成立。
            \item \textbf{牛顿第三定律}：两物体间的作用力与反作用力大小相等、方向相反、沿同一直线、分别作用在两个物体上、性质相同、同时存在同时消失。由于不涉及运动，它对参考系没有限制。
        \end{enumerate}
        惯性系的确认靠实验。以太阳中心为原点、指向任一恒星的坐标系是较好的惯性系；地球有自转和公转，法向加速度很小（$7.3\times10^{-5}$ 与 $2.0\times10^{-7}\ \mathrm{rad\cdot s^{-1}}$），在一般精度内可把地面参考系近似看作惯性系。

        \subsubsection{自然界的基本相互作用与常见的力}
        按性质，自然界的一切相互作用可归结为四种：强相互作用、电磁相互作用、弱相互作用、引力相互作用。日常遇到的摩擦力、弹性力、张力、支承力、浮力、黏性力等，都是分子（原子）间电磁相互作用的宏观表现。

        \begin{center}
        \begin{tabular}{lll}
        \toprule
        力 & 表达式 & 说明 \\
        \midrule
        重力 & $\boldsymbol{G}=m\boldsymbol{g}$ & $\boldsymbol{g}$ 为重力加速度，竖直向下 \\
        弹力（弹簧） & $F=-kx$ & 胡克定律，$x$ 为相对原长的形变 \\
        滑动摩擦力 & $F_{\mathrm{f}}=\mu F_{\mathrm{N}}$ & 方向与相对滑动方向相反 \\
        静摩擦力 & $0\leqslant F_{\mathrm{f}}\leqslant\mu_{\mathrm{s}}F_{\mathrm{N}}$ & 由平衡条件确定 \\
        万有引力 & $F=G\dfrac{m_{1}m_{2}}{r^{2}}$ & $G=6.672\times10^{-11}\ \mathrm{N\cdot m^{2}\cdot kg^{-2}}$ \\
        \bottomrule
        \end{tabular}
        \end{center}

        \subsubsection{非惯性系与惯性力}
        在相对惯性系以加速度 $\boldsymbol{a}_{0}$ 平动的参考系中，牛顿定律不成立。若仍要套用牛顿第二定律的形式，须引入\textbf{惯性力}
        \[
        \boldsymbol{F}_{0}=-m\boldsymbol{a}_{0}
        \]
        于是有 $\boldsymbol{F}+\boldsymbol{F}_{0}=m\boldsymbol{a}'$。惯性力方向与非惯性系相对惯性系的加速度方向相反，大小等于研究对象的质量乘该加速度。惯性力无施力物体、无反作用力，它不是物体间的相互作用，只是物体惯性在非惯性系中的表现。
        在匀速转动的参考系（角速度 $\omega$，质点到轴距离 $r$）中，相对该系静止的物体还受\textbf{惯性离心力}
        \[
        \boldsymbol{F}_{0}=m\omega^{2}\boldsymbol{r}
        \]
        方向沿半径向外。\corepoint{惯性离心力不是向心力的反作用力：它作用在小球上，而向心力的反作用力（常称离心力）作用在绳子上。}

    \subsection{动量定理\texorpdfstring{\hspace{0.5em}}{ }动量守恒定律}
        \subsubsection{冲量与质点的动量定理}
        力对时间的累积量称为\textbf{冲量}（impulse）：
        \[
        \boldsymbol{I}=\int_{t_{0}}^{t}\boldsymbol{F}\,\mathrm{d}t,\qquad \mathrm{d}\boldsymbol{I}=\boldsymbol{F}\,\mathrm{d}t
        \]
        由 $\boldsymbol{F}=\mathrm{d}\boldsymbol{p}/\mathrm{d}t$ 得积分形式的\textbf{质点动量定理}：
        \keyformula{\boldsymbol{I}=\int_{t_{0}}^{t}\boldsymbol{F}\,\mathrm{d}t=m\boldsymbol{v}-m\boldsymbol{v}_{0}=\boldsymbol{p}-\boldsymbol{p}_{0}}
        即质点所受合外力的冲量等于质点动量的增量。冲量是矢量，方向与动量增量相同；单位 $\mathrm{N\cdot s}$，量纲 $\mathrm{MLT^{-1}}$。分量形式为
        \[
        I_{x}=\int F_{x}\,\mathrm{d}t=mv_{x}-mv_{x0},\quad
        I_{y}=\int F_{y}\,\mathrm{d}t=mv_{y}-mv_{y0},\quad
        I_{z}=\int F_{z}\,\mathrm{d}t=mv_{z}-mv_{z0}
        \]
        碰撞、爆炸等过程中冲力变化复杂、作用时间极短，可用\textbf{平均冲力}代替瞬时冲力：
        \[
        \bar{\boldsymbol{F}}\Delta t=m\boldsymbol{v}-m\boldsymbol{v}_{0}
        \]

        \subsubsection{质点系的动量定理}
        系统内力总是成对出现且矢量和为零，内力的冲量只改变系统内动量的分配，不改变系统的总动量，因此
        \keyformula{\boldsymbol{F}_{\text{外}}=\frac{\mathrm{d}\boldsymbol{p}}{\mathrm{d}t},\qquad
        \int_{t_{0}}^{t}\boldsymbol{F}_{\text{外}}\,\mathrm{d}t=\boldsymbol{p}-\boldsymbol{p}_{0}=\Delta\boldsymbol{p}}

        \subsubsection{动量守恒定律}
        若系统所受合外力为零，则系统总动量保持不变：
        \keyformula{\boldsymbol{p}=\sum_{i}m_{i}\boldsymbol{v}_{i}=\text{常矢量}}
        应用时应注意：
        \begin{enumerate}
            \item “合外力为零”指外力的矢量和为零。此时系统总动量守恒，但系统内各质点的动量可以转移和传递；
            \item 若合外力不为零，但在某一方向上外力分量之和为零，则该方向上的动量分量守恒：$F_{x}=0$ 时 $p_{x}=\sum m_{i}v_{ix}=$ 常量，其余方向同理；
            \item 碰撞、爆炸等问题中冲力极大、作用时间极短，一般外力（如重力）可忽略，可近似按动量守恒处理；
            \item 动量守恒定律不仅适用于宏观物体，也适用于原子、分子、光子等微观粒子间的相互作用，是自然界最普遍、最基本的定律之一。
        \end{enumerate}

    \subsection{质心\texorpdfstring{\hspace{0.5em}}{ }质心运动定理}
        \subsubsection{质心位置}
        质点系的\textbf{质心}（center of mass）是质量加权平均的位置，其位矢为
        \keyformula{\boldsymbol{r}_{C}=\frac{\sum_{i}m_{i}\boldsymbol{r}_{i}}{\sum_{i}m_{i}}=\frac{1}{m}\sum_{i}m_{i}\boldsymbol{r}_{i},\quad m=\sum_{i}m_{i}}
        质量连续分布的物体用积分代替求和：
        \keyformula{\boldsymbol{r}_{C}=\frac{1}{m}\int_{V}\boldsymbol{r}\,\mathrm{d}m,\qquad
        x_{C}=\frac{1}{m}\int x\,\mathrm{d}m,\quad
        y_{C}=\frac{1}{m}\int y\,\mathrm{d}m,\quad
        z_{C}=\frac{1}{m}\int z\,\mathrm{d}m}

        \subsubsection{质心速度与质心加速度}
        \[
        \boldsymbol{v}_{C}=\frac{\mathrm{d}\boldsymbol{r}_{C}}{\mathrm{d}t}=\frac{1}{m}\sum_{i}m_{i}\boldsymbol{v}_{i},
        \qquad
        \boldsymbol{a}_{C}=\frac{\mathrm{d}\boldsymbol{v}_{C}}{\mathrm{d}t}=\frac{1}{m}\sum_{i}m_{i}\boldsymbol{a}_{i}
        \]

        \subsubsection{质心运动定理}
        由系统动量 $\boldsymbol{p}=m\boldsymbol{v}_{C}$ 及质点系动量定理得
        \keyformula{\boldsymbol{F}_{\text{外}}=m\boldsymbol{a}_{C}=m\frac{\mathrm{d}\boldsymbol{v}_{C}}{\mathrm{d}t}}
        系统合外力等于总质量乘质心加速度，即质心的运动与一个质量等于系统总质量、受同样合外力的质点的运动完全相同。内力不改变质心的运动状态；合外力为零时，质心保持静止或匀速直线运动。炮弹在空中爆炸时，弹片飞向各方，但质心仍沿原抛物线运动。

    \subsection*{火箭飞行原理（变质量系统，选读）}
    火箭是靠排出内部燃烧室产生的气体获得推力的。设某时刻火箭质量为 $m$、速度为 $\boldsymbol{v}$，取火箭与燃料为研究对象，在 $\mathrm{d}t$ 时间内喷出质量为 $-\mathrm{d}m$（$\mathrm{d}m<0$）、相对火箭速度为 $\boldsymbol{u}$ 的气体，由质点系动量定理可得\textbf{火箭方程}
    \[
    m\frac{\mathrm{d}\boldsymbol{v}}{\mathrm{d}t}=\boldsymbol{F}+\boldsymbol{u}\frac{\mathrm{d}m}{\mathrm{d}t}
    \]
    其中 $\boldsymbol{u}\,\mathrm{d}m/\mathrm{d}t$ 称为火箭发动机的\textbf{推力}。当外力 $\boldsymbol{F}=0$（远离大气层或在太空）时积分得\textbf{齐奥尔科夫斯基公式}
    \keyformula{v=v_{0}+u\ln\frac{m_{0}}{m_{\text{火箭}}}}
    火箭增加的速度不仅与喷气速度 $u$ 有关，还与始末质量比 $N=m_{0}/m_{\text{火箭}}$ 有关，质量比越大末速度越大。单级火箭受质量比限制，末速度达不到第一宇宙速度，故发射航天器需用多级火箭。

\section{角动量\texorpdfstring{\hspace{0.5em}}{ }角动量守恒定律}
    \subsection{质点的角动量}
    设质量为 $m$ 的质点，某时刻对惯性系中参考点 $O$ 的位矢为 $\boldsymbol{r}$，动量为 $\boldsymbol{p}=m\boldsymbol{v}$，则质点对该点的\textbf{角动量}（angular momentum，又称动量矩）为
    \keyformula{\boldsymbol{L}=\boldsymbol{r}\times\boldsymbol{p}=\boldsymbol{r}\times m\boldsymbol{v}}
    大小为
    \[
    L=rp\sin\varphi=mrv\sin\varphi
    \]
    其中 $\varphi$ 是 $\boldsymbol{r}$ 与 $\boldsymbol{v}$ 的夹角。$\boldsymbol{L}$ 的方向垂直于 $\boldsymbol{r}$ 与 $\boldsymbol{v}$ 所决定的平面，指向由右手螺旋定则确定；单位 $\mathrm{kg\cdot m^{2}\cdot s^{-1}}$，量纲 $\mathrm{ML^{2}T^{-1}}$。
    角动量与所选参考点有关，说明角动量时必须指明是对哪一个固定点而言的。定义本身并不要求质点作圆周运动：作匀速直线运动的质点，对某固定点的角动量也是恒定的。
    质点系的总角动量为各质点角动量的矢量和：
    \[
    \boldsymbol{L}=\sum_{i}\boldsymbol{L}_{i}=\sum_{i}\boldsymbol{r}_{i}\times m_{i}\boldsymbol{v}_{i}
    \]

    \subsection{质点角动量的时间变化率\texorpdfstring{\hspace{0.5em}}{ }力矩}
    作用于质点的力 $\boldsymbol{F}$ 对参考点 $O$ 的\textbf{力矩}（moment of force）定义为
    \keyformula{\boldsymbol{M}=\boldsymbol{r}\times\boldsymbol{F},\qquad M=rF\sin\varphi=Fd}
    其中 $d=r\sin\varphi$ 为力臂。力矩是矢量，方向垂直于 $\boldsymbol{r}$ 与 $\boldsymbol{F}$ 所决定的平面，由右手螺旋定则确定；单位 $\mathrm{N\cdot m}$。
    对质点系，一对内力对同一参考点的力矩矢量和恒为零（内力沿两质点连线，与连线共线），故系统所受合外力矩等于各外力矩的矢量和：
    \[
    \boldsymbol{M}=\sum_{i}\boldsymbol{r}_{i}\times\boldsymbol{F}_{i}^{\text{外}}
    \]

    \subsection{质点的角动量定理\texorpdfstring{\hspace{0.5em}}{ }角动量守恒定律}
        \subsubsection{角动量定理}
        对 $\boldsymbol{L}=\boldsymbol{r}\times m\boldsymbol{v}$ 求导，并注意 $\boldsymbol{v}\times m\boldsymbol{v}=0$，得角动量定理的微分形式
        \keyformula{\boldsymbol{M}=\frac{\mathrm{d}\boldsymbol{L}}{\mathrm{d}t}}
        即质点（质点系）所受对某参考点的合外力矩，等于同一参考点角动量对时间的变化率。对时间积分得积分形式（冲量矩定理）：
        \keyformula{\int_{t_{1}}^{t_{2}}\boldsymbol{M}\,\mathrm{d}t=\boldsymbol{L}_{2}-\boldsymbol{L}_{1}}
        $\displaystyle\int_{t_{1}}^{t_{2}}\boldsymbol{M}\,\mathrm{d}t$ 称为\textbf{冲量矩}（moment of impulse），是外力矩对时间的累积量，它等于系统角动量的增量。

        \subsubsection{角动量守恒定律}
        若对某固定参考点合外力矩为零（$\boldsymbol{M}=0$），则系统对该点的角动量保持不变：
        \keyformula{\boldsymbol{L}=\text{常矢量}}
        这就是\textbf{角动量守恒定律}。应用要点：
        \begin{enumerate}
            \item 有心力（力的作用线始终通过某定点——力心）对力心的力矩恒为零，只受有心力作用的系统角动量守恒；
            \item 合外力矩虽不为零，但对某轴的分量为零时，该轴上的角动量分量守恒；
            \item 角动量守恒与动量守恒、能量守恒一样，是自然界的普遍规律，在微观领域（粒子衰变、原子物理）同样适用。
        \end{enumerate}
        \corepoint{判断角动量是否守恒只看对所选点的合外力矩，不看合外力；行星绕日运动时动量不守恒，但角动量守恒。}

        \subsubsection{应用：开普勒第二定律（面积定律）}
        行星只受太阳的有心力作用，对太阳的角动量守恒。设极短时间 $\Delta t$ 内矢径扫过面积 $\Delta S=\dfrac{1}{2}r|\Delta\boldsymbol{r}|\sin\varphi$，则
        \[
        L=rmv\sin\varphi=2m\lim_{\Delta t\to0}\frac{\Delta S}{\Delta t}=2m\frac{\mathrm{d}S}{\mathrm{d}t}
        \]
        故\textbf{面积速率}为常量：
        \keyformula{\frac{\mathrm{d}S}{\mathrm{d}t}=\frac{L}{2m}=\text{常量}}
        即行星对太阳的矢径在相等时间内扫过相等的面积。

\section{能量\texorpdfstring{\hspace{0.5em}}{ }能量守恒定律}
    \subsection{动能\texorpdfstring{\hspace{0.5em}}{ }功\texorpdfstring{\hspace{0.5em}}{ }动能定理}
        \subsubsection{动能}
        \keyformula{E_{\mathrm{k}}=\frac{1}{2}mv^{2}}
        动能是标量，单位焦耳（$\mathrm{J}$），量纲 $\mathrm{ML^{2}T^{-2}}$。质点系的总动能是各质点动能的代数和：
        \[
        E_{\mathrm{k}}=\sum_{i}\frac{1}{2}m_{i}v_{i}^{2}
        \]

        \subsubsection{功与功率}
        力在位移方向上的分量与位移大小的乘积称为\textbf{功}（work）。元功为
        \keyformula{\mathrm{d}W=\boldsymbol{F}\cdot\mathrm{d}\boldsymbol{r}=F\cos\alpha\,\mathrm{d}r}
        沿路径 $A\to B$ 的总功为线积分：
        \keyformula{W=\int_{A}^{B}\boldsymbol{F}\cdot\mathrm{d}\boldsymbol{r}
        =\int_{A}^{B}\left(F_{x}\mathrm{d}x+F_{y}\mathrm{d}y+F_{z}\mathrm{d}z\right)}
        功是标量，有正负：$\alpha<\pi/2$ 时做正功，$\alpha>\pi/2$ 时做负功，$\alpha=\pi/2$ 时不做功。功的四个性质：
        \begin{enumerate}
            \item 合力的功等于各分力功的代数和；
            \item 功是\textbf{过程量}，依赖于路径，某一时刻没有“功”；
            \item 功是\textbf{相对量}，位移与参考系有关，同一力在不同参考系中做功不同；
            \item 一对相互作用力的总功由相对位移决定，无相对位移时内力总功为零。
        \end{enumerate}
        瞬时功率为
        \keyformula{P=\frac{\mathrm{d}W}{\mathrm{d}t}=\boldsymbol{F}\cdot\boldsymbol{v}}
        单位瓦特（$\mathrm{W}$）。

        \subsubsection{质点与质点系的动能定理}
        \keyformula{W=\int_{A}^{B}\boldsymbol{F}\cdot\mathrm{d}\boldsymbol{r}
        =\frac{1}{2}mv_{B}^{2}-\frac{1}{2}mv_{A}^{2}=E_{\mathrm{k}B}-E_{\mathrm{k}A}}
        即合外力对质点做的总功等于质点动能的增量。对质点系，
        \[
        W_{\text{外}}+W_{\text{内}}=\sum_{i}\frac{1}{2}m_{i}v_{i}^{2}-\sum_{i}\frac{1}{2}m_{i}v_{i0}^{2}
        \]
        \corepoint{动能定理只适用于惯性系；功是过程量、动能是状态量，但动能定理使“求速率—求位置”类问题不必追究中间过程。}

    \subsection{保守力\texorpdfstring{\hspace{0.5em}}{ }势能\texorpdfstring{\hspace{0.5em}}{ }功能原理}
        \subsubsection{保守力与非保守力}
        做功只由始末位置决定、与路径无关的力称为\textbf{保守力}（conservative force），等价条件是沿任意闭合路径做功为零：
        \keyformula{\oint_{L}\boldsymbol{F}_{\text{保}}\cdot\mathrm{d}\boldsymbol{r}=0}
        重力、弹性力、万有引力、静电力都是保守力；摩擦力、黏性力、流体阻力、爆炸力做功与路径有关，称为\textbf{非保守力}（耗散力）。

        \subsubsection{三种保守力的功与势能}
        \begin{center}
        \small
        \setlength{\tabcolsep}{9pt}
        \begin{tabular}{lll}
        \toprule
        保守力 & 做功 & 势能（势能零点） \\
        \midrule
        重力 $mg$ & $W=-mg\left(y_{b}-y_{a}\right)$ & $E_{\mathrm{p}}=mgh$（地面） \\
        弹性力 $-kx$ & $W=-\left(\frac{1}{2}kx_{b}^{2}-\frac{1}{2}kx_{a}^{2}\right)$
        & $E_{\mathrm{p}}=\frac{1}{2}kx^{2}$（原长处） \\
        万有引力 & $W=-Gm_{1}m_{2}\left(\frac{1}{r_{b}}-\frac{1}{r_{a}}\right)$
        & $E_{\mathrm{p}}=-G\frac{m_{1}m_{2}}{r}$（无穷远） \\
        \bottomrule
        \end{tabular}
        \end{center}
        保守力做功等于势能增量的负值：
        \keyformula{W_{\text{保}}=-\Delta E_{\mathrm{p}}=-\left(E_{\mathrm{p}b}-E_{\mathrm{p}a}\right)}
        由此可得保守力与势能的微分关系（梯度关系）：
        \keyformula{\boldsymbol{F}_{\text{保}}=-\nabla E_{\mathrm{p}}
        =-\left(\frac{\partial E_{\mathrm{p}}}{\partial x}\boldsymbol{i}
        +\frac{\partial E_{\mathrm{p}}}{\partial y}\boldsymbol{j}
        +\frac{\partial E_{\mathrm{p}}}{\partial z}\boldsymbol{k}\right)}
        \corepoint{势能属于相互作用系统所共有，单个物体无“势能”；且只有存在保守力时才能引入势能。}

        \subsubsection{功能原理}
        把内力分为保守内力与非保守内力，并用势能增量代替保守内力的功，得质点系的\textbf{功能原理}：
        \keyformula{W_{\text{外}}+W_{\text{非保内}}=E-E_{0},\qquad E=E_{\mathrm{k}}+E_{\mathrm{p}}}
        即系统所有外力做的功与非保守内力做的功之和，等于系统机械能的增量。功能原理与质点系动能定理物理内容相同，区别只在于把保守内力的功写成了势能增量的负值。

    \subsection{机械能守恒定律与能量守恒定律}
        \subsubsection{机械能守恒定律}
        若 $W_{\text{外}}=0$ 且 $W_{\text{非保内}}=0$（只有保守内力做功），则系统机械能保持不变：
        \keyformula{E=E_{\mathrm{k}}+E_{\mathrm{p}}=\text{常量}}
        机械能守恒定律只适用于惯性系，且与参考系的选择有关（在非惯性系中惯性力可能做功）。

        \subsubsection{碰撞}
        碰撞是伴随形变与形变恢复的短暂过程，系统动量总是守恒的：
        \begin{enumerate}
            \item \textbf{弹性碰撞}：形变完全恢复，机械能（动能）守恒；
            \item \textbf{完全非弹性碰撞}：形变完全不能恢复，碰后两物体一起运动，动能损失最大；
            \item \textbf{非弹性碰撞}：碰后保留部分形变，机械能有损失。
        \end{enumerate}

        \subsubsection{普遍的能量守恒定律}
        与外界无能量交换的系统称为孤立系统。在孤立系统内，各种形式的能量可以相互转化，但系统总能量保持不变：能量既不能创造，也不能消灭，只能从一种形式转化为另一种形式。这就是\textbf{能量转化和守恒定律}，是自然界最普遍的规律。

        \subsubsection{应用：宇宙速度与黑洞模型}
        取无穷远为引力势能零点，物体从地球表面（半径 $R$）以速率 $v$ 逃逸到无穷远且末速为零，由机械能守恒
        \[
        \frac{1}{2}mv^{2}-G\frac{mM}{R}=0
        \]
        并利用 $G\dfrac{M}{R^{2}}=g$，得\textbf{第二宇宙速度（逃逸速度）}
        \keyformula{v=\sqrt{2gR}=\sqrt{\frac{2Gm_{\text{天体}}}{R}}}
        代入地球数据得 $v\approx11.2\ \mathrm{km\cdot s^{-1}}$。若某天体密度极大而半径极小，使逃逸速度接近光速，则连它发出的光也无法逃离，这种天体称为\textbf{黑洞}（black hole）。
